%-------------------------
% Cover Letter - Radix Trading Quantitative Technologist (C++)
% Author: Ajay Venkatesh
%-------------------------

\documentclass[letterpaper,11pt]{article}

\usepackage[top=0.8in, bottom=0.8in, left=0.9in, right=0.9in]{geometry}
\usepackage[hidelinks]{hyperref}
\usepackage{lmodern}
\usepackage{parskip}
\usepackage{microtype}

\setlength{\parskip}{6pt}
\setlength{\parindent}{0pt}

\begin{document}

\begin{flushleft}
\textbf{\large Ajay Venkatesh} \\
Brooklyn, NY \\
+1 516 585 9013 \\
\href{mailto:av3855@nyu.edu}{av3855@nyu.edu}
\end{flushleft}

\vspace{10pt}

May 10, 2026

\vspace{6pt}

Hiring Manager \\
Radix Trading LLC

\vspace{10pt}

Dear Hiring Manager,

My name is Ajay Venkatesh --- finishing my M.S. in Computer Engineering at NYU, with several years of production software experience: a few years at PwC in Bengaluru, a summer SDE internship at Amazon, and ongoing on-campus work at NYU's Novel AI Technologies. My background sits at the intersection of \textbf{systems-level engineering, large-scale data processing,} and \textbf{applied ML} --- the same triangle the Quantitative Technologist role lives in.

The Radix posting maps onto a few of my strongest threads. My graduate coursework was deliberately weighted toward the low end of the stack --- \textbf{Computer Architecture, System Design, Data Structures, Big Data} --- alongside \textbf{Machine Learning} and \textbf{Deep Learning}. \textbf{C++} is where I cut my teeth and where I keep going back when the abstractions stop being honest about what's happening underneath. I've enjoyed implementing things from scratch over leaning on libraries when the problem rewards understanding the nuance of the language and the machine.

On data scale, my Big Data project built a \textbf{PySpark / HDFS} pipeline over millions of flight records, training Random Forest and Gradient Boosted Tree models with rigorous EDA --- small relative to a quant firm's petabyte-scale data, but the same shape: feature pipelines, anomaly detection, and squeezing accuracy out of irregular real-world inputs. At Amazon I deployed a \textbf{Retrieval-Augmented Generation} pipeline with an \textbf{agentic workflow} for incident reasoning, implemented a \textbf{Model Context Protocol (MCP)} server that materially reduced query time and compute costs, and built a high-throughput, event-driven triage system on \textbf{ECS Fargate, SQS, and SNS} --- the kind of distributed, low-latency systems work that translates to live-trading infrastructure.

What pulls me toward Radix is the combination the posting describes: a hacker's approach to ideas, modern C++ pushed daily, ML research and simulation infrastructure over enormous data, and a research-driven culture that doesn't romanticize the academic version of any of it. I want to work somewhere the work matters in real units --- profit, latency, throughput --- and where the team is sharper than I am so I can keep learning at speed. The Radix description reads like that environment.

I'd welcome the chance to discuss further. Thanks for considering my application.

\vspace{12pt}

Sincerely, \\
Ajay Venkatesh

\end{document}
